% Classe per tesi di Ingegneria Informatica all'Universita Roma Tre
\documentclass{TesiDiaUniroma3}

\usepackage[utf8]{inputenc}
\usepackage[T1]{fontenc}
\usepackage[italian]{babel}
\usepackage{lmodern}
\usepackage{microtype}
\usepackage{amsmath}
\usepackage{amssymb}
\usepackage{booktabs}
\usepackage{tabularx}
\usepackage{array}
\usepackage[table]{xcolor}
\usepackage{tikz}
\usetikzlibrary{arrows.meta,positioning,shapes.geometric,fit,calc,backgrounds}
\usepackage{url}
\usepackage{xurl}
\usepackage[plainpages=false,hidelinks]{hyperref}
\usepackage[all]{hypcap}
\hypersetup{
  pdftitle={Agent-based Earth Observation discovery: un approccio multi-agentico per l'accesso al dato nei cataloghi di osservazione della Terra},
  pdfauthor={Alessio Pecilli},
  pdfsubject={Architettura e valutazione di un workspace conversazionale per fonti geospaziali},
  pdfkeywords={Copernicus, osservazione della Terra, agenti LLM, MCP, STAC},
  pdfdisplaydoctitle=true
}

\bibliographystyle{plain}

% Dati del frontespizio.
\titolo{Agent-based Earth Observation discovery: un approccio multi-agentico per
l'accesso al dato nei cataloghi di osservazione della Terra}
\autore{Alessio Pecilli}
\matricola{589232}
\relatore{Prof. Matteo Rosati}
\annoAccademico{2025/2026}

\begin{document}

\frontmatter
\generaFrontespizio
% Da completare in seguito.

\chapter*{Sommario}
\addcontentsline{toc}{chapter}{Sommario}

La consultazione delle fonti di osservazione della Terra richiede di distinguere
eventi naturali, acquisizioni satellitari, attivazioni cartografiche e cataloghi
di dataset. Questa tesi presenta Copernicus Data Workspace, un prototipo
conversazionale che interpreta richieste geografiche e coordina il recupero
di informazioni mantenendone riconoscibili provenienza e limiti.

Il sistema adotta un workflow LangGraph con componenti decisionali basati su
modelli linguistici e trasformazioni deterministiche. Il percorso di Discovery
interroga un catalogo di metadati NebulaGraph alimentato da harvester Marine,
STAC e Drought. Il percorso operativo consulta GDACS, Copernicus Emergency
Management Service e il catalogo STAC remoto; MCP delimita le integrazioni
GDACS e CEMS. Memoria conversazionale, artefatti e telemetria sostengono
rispettivamente i turni successivi, la presentazione dei risultati e la diagnosi.
Conteggi e aggregazioni delle visualizzazioni derivano dai record recuperati,
senza attribuire al modello il ruolo di sorgente dei valori.

La valutazione distingue controlli software, benchmark del workflow e verifica
dell'interfaccia. Il protocollo descritto combina controlli deterministici e
LLM-as-a-Judge, rendendo espliciti denominatori, dipendenze tra casi e limiti
del valutatore. Sul catalogo a grafo il sistema supera circa due casi su tre,
con il limite principale nell'ordinamento dei risultati; sulle acquisizioni
satellitari supera tutti i casi del benchmark. Il contributo è progettuale e implementativo:
non viene dimostrata una superiorità quantitativa rispetto a un agente singolo
o a una pipeline più semplice. Restano necessari confronti controllati,
validazione umana dei giudizi e prove conversazionali nel browser.

\clearpage

\generaIndice

\mainmatter
\capitolo{Introduzione}{introduzione}
\capitolo{Dall'intelligenza artificiale agli agenti basati su LLM}{1}
\capitolo{Model Context Protocol: contesto, strumenti e interoperabilità}{2}
\capitolo{Fonti informative geospaziali: GDACS e Copernicus}{3}
\capitolo{Architettura del sistema proposto}{4}
\capitolo{Funzionamento end-to-end, casi d'uso e verifica}{5}
\capitolo{Osservabilità del workflow e diagnosi}{6}
\capitolo{Valutazione con LLM-as-a-Judge}{7}
\conclusioni{conclusioni}

\backmatter
\begingroup
\small
\linespread{1.2}\selectfont
\bibliography{bibliografia}
\endgroup

\end{document}
